\section{Dialectics}

\begin{quotex}
Wer sagt A, muss B sagen.
\end{quotex}


This is the dictionary definition of dialectic:

\begin{quotex}
The art or practice of arriving at the truth by the exchange of logical arguments.
\end{quotex}


In a Socratic-style dialog, a

\begin{quotex}
series of questions clarifies a more precise statement of a vague belief, logical consequences of that statement are explored, and a contradiction is discovered (Wikipedia)
\end{quotex}


As a general principle, therefore, in regards to to comment, call it ``A'', there can be two responses ``B''.

\begin{enumerate}


\item Expose the unexpressed presuppositions behind the comment

\item Follow up on the logical consequences of the comment
\end{enumerate}


So if someone says A, he is intellectually obliged to accept B. Then you can answer B with C so that a legitimate dialogue may ensue. If you simply repeat A, usually on the assumption that I am incapable of grasping your original comment, then you have ended the dialog. Your future comments will be blocked.

Don't presume you understand my inner psychological states. If you are incapable of understanding dialectic, then I might respond with rhetoric. It's just part of the game; learn to play Calcio fiorentino, not Tiddlywinks.

The number of readers who take this blog seriously are few; they don't want to be bothered with nonsense. Some of them do take the trouble to post relevant and useful comments, but most of them just reach out to me privately.

A comment should reference the text and not bring in writers or concepts not referenced. E.g., if I never mentioned ``intersubjectivity'', then you shouldn't either.

No one has a right to waste my time reading poorly thought-out comments. There is a web site that charges \$120/year so that your comments will make it to the top of the list, tout suite. I wouldn't think that there would be enough suckers to fall for that, but I've been mistaken before. Perhaps I'll try it. Just name a fair price to let your comments flow through unmoderated.

Ultimately, this is for your own good:

\begin{quotex}
The principal aim of Socratic activity may be to improve the soul of the interlocutors, by freeing them from unrecognized errors; or indeed, by teaching them the spirit of inquiry. (Wikipedia)
\end{quotex}



\flrightit{Posted on 2021-07-10 by Cologero}

